\chapter{Chiral Matter and Internal Symmetries}
\label{ch:ch06_chiral}

The question of this chapter is which chiral content the tower carries at three constituents and which of its symmetries can be
gauged. It treats the projectors of rank $64$, among them the family projector $\Pfam$, which is kept apart from the chiral projector;
the chiral Pati--Salam block $P_*$, which carries one family with a right-handed neutrino; and the charges and anomalies of the block. The
content is proved. The choice of $P_*$ among the projectors is constructed, not selected, and chirality is imposed, not derived
(\cref{prop:chiralpreserved}).

% --- source: tex/s07_projectors#0 ((head))


% --- source: tex/s07_projectors#1 (Chirality, projectors, and colour)
\section{Chirality, projectors, and colour}
\label{sec:projectors}

% --- source: tex/s07_projectors#2 (Unprojected matter is vector-like)
\subsection{Unprojected matter is vector-like}

With $i=J$ and $\gamma_5=\Gamma$ as in \cref{thm:i}, consider the part of the Lorentz commutant that preserves chirality. It splits
into a \emph{vector class} $\Cl^0(W)$, the even subalgebra of the internal directions, whose elements commute with the interior gamma
matrices, and an \emph{axial class} $\Omega\,\Cl^0(W)$, whose elements anticommute with them.

\begin{proposition}[Unprojected matter is vector-like]
\label{prop:vectorlike}
Every grade-two internal algebra lies in the vector class. For a gauge algebra in the vector class, $\gamma_\mu$ is an invertible,
complex-linear, gauge-equivariant map between the two chirality sectors, so the unprojected matter is vector-like at every rung.
\end{proposition}

\begin{proposition}
\label{prop:axialtriple}
The two classes form a symmetric pair, $[\mathrm a,\mathrm a]\subseteq\mathrm v$, so the axial class is a Lie triple system and no
nonabelian gauge algebra lies entirely inside it.
\end{proposition}

This does not exclude generators with both parts. If $[\Gamma,t_a]=0$ and $P_L=\tfrac12(\id-\Gamma)$, then
\[
[P_Lt_a,P_Lt_b]=P_L^2[t_a,t_b]=\epsilon_{abc}P_Lt_c ,
\]
a nonabelian algebra, half vector and half axial. Such \emph{dressing} is available to nonabelian carriers exactly as to abelian ones;
only the purely axial summand is excluded. For scalar dressings $f(\Gamma)t_a$ with $f=a+b\Gamma$ and unchanged structure constants,
closure forces $f^2=f$, that is $a^2+b^2=a$ and $2ab=b$, so $f\in\{0,\id,P_L,P_R\}$.

% --- source: tex/s07_projectors#3 (A family of projectors at three constituents)
\subsection{A family of projectors at three constituents}

For the undressed internal bivector action of \cref{prop:vectorlike,prop:axialtriple}, chiral matter requires a projection; with dressing and
spectator states, \cref{thm:p1chiral} obtains a chiral interaction otherwise. At $n=3$ write $J^{(k)}=e^{(k)}_2e^{(k)}_3$ for the surplus bivector and
$\omega^{(k)}=e^{(k)}_1\cdots e^{(k)}_6$ for the volume of constituent $k$, and $K^{(3)}_{14}=e^{(3)}_1e^{(3)}_4$ for the pole bivector
joining the first temporal and the first spatial direction of constituent three. The three operators
\[
X_1=J^{(1)}J^{(2)},\qquad X_2=\omega^{(3)},\qquad X_3=\omega^{(1)}K^{(3)}_{14}
\]
are commuting symmetric traceless involutions, and they generate the orthogonal projectors $P_1=\tfrac12(\id+X_1)$ of rank $256$,
$P_{ij}=\tfrac14(\id+X_i)(\id+X_j)$ of rank $128$, and $\Pfam=\tfrac18\prod_i(\id+X_i)$ of rank $64$.

\begin{definition}[Ambient and faithful algebra]
\label{def:keff}
For an orthogonal projector $P$ commuting with the Lorentz action let $\mathfrak g_{\rm phys}(P)$ be the internal bivectors commuting
with $P$, $\mathfrak k(P)$ its compact part (the $+1$ eigenspace of $\theta(X)=-X^{\!\top}$), and
\[
\keff(P)=\mathfrak k(P)\big/\{X\in\mathfrak k(P):XP=0\}.
\]
\end{definition}

The quotient is not a technicality: an element can preserve the block and annihilate every state in it.

\begin{example}[An annihilator]
\label{ex:annihilator}
With $X_1=J^{(1)}J^{(2)}$, $(J^{(1)})^2=(J^{(2)})^2=-\id$ and $[J^{(1)},J^{(2)}]=0$,
\[
(J^{(1)}+J^{(2)})(\id+X_1)=J^{(1)}+J^{(2)}+(J^{(1)})^2J^{(2)}+J^{(2)}J^{(1)}J^{(2)}=J^{(1)}+J^{(2)}-J^{(2)}-J^{(1)}=0 .
\]
So $J^{(1)}+J^{(2)}$ is a compact element of $\mathfrak k(\Pfam)$ that annihilates the block of $\Pfam$, although it does not annihilate the
block of $P_{23}$. A table of gauge algebras that omits the quotient overcounts, here by three dimensions.
\end{example}

\begin{theorem}[The projector family]
\label{thm:table}
Computed on the $512$-dimensional module as the exact nullspace of $X\mapsto[X,P]$ over the $\binom{14}2=91$ internal bivectors, with
$\keff$ decomposed on the image ($\dim\mathfrak g_{\rm phys}$, $\dim\mathfrak k$ and $\dim\mathfrak p$ count the commuting algebra before the quotient
by its kernel on the block, so $\dim\keff=\dim\mathfrak k-\dim\ker$):
\begin{center}
\begin{tabular}{@{}lrrrrrlc@{}}
\toprule
$P$ & rank & $\dim\mathfrak g_{\rm phys}$ & $\dim\mathfrak k$ & $\dim\mathfrak p$ & $\dim\ker$ & $\keff$ & axial\\
\midrule
$\Pfam$ & $64$ & $17$ & $9$ & $8$ & $3$ & $\uu1^3\oplus\su2$ & yes\\
$P_{23}$ & $128$ & $23$ & $9$ & $14$ & $0$ & $\uu1^3\oplus\su2\oplus\su2$ & yes\\
$P_{13}$ & $128$ & $33$ & $17$ & $16$ & $3$ & $\uu1\oplus\su2\oplus\so5$ & yes\\
$P_{12}$ & $128$ & $27$ & $15$ & $12$ & $3$ & $\su2^{\oplus4}$ & no\\
$P_1$ & $256$ & $51$ & $27$ & $24$ & $3$ & $\so6\oplus\su2^{\oplus3}$ & no\\
\bottomrule
\end{tabular}
\end{center}
Here a projector is axial when $P\gamma^\mu P=0$. Factors are identified by dimension and rank: $(3,1)=\su2$, $(10,2)=\so5\cong\spp2$,
$(15,3)=\so6\cong\su4$; a six-dimensional semisimple compact factor is necessarily $\su2\oplus\su2$. The centres have dimensions
$3,3,1,0,0$ and the derived algebras $3,6,13,12,24$.
\end{theorem}

\noindent Executable check: a4-4 and a10-2 (exact; \cref{app:methods}).

\begin{corollary}[No colour on the axial members]
\label{cor:nocolour}
The largest simple factor among the rank-$64$ and rank-$128$ blocks is $\so5$, every proper subalgebra of which has dimension at most
six. Since $\su3$ is simple of dimension eight, it embeds in none of them. Dimension eight is necessary and not sufficient: $P_{13}$
carries an $\so5$ of dimension ten and still contains no $\su3$, so the test is an embedding test.
\end{corollary}

% --- source: tex/s07_projectors#4 (Colour at rank 256)
\subsection{Colour at rank 256}

\begin{theorem}[Colour on $P_1$]
\label{thm:su3}
Let $f_1,\dots,f_6$ be the six spatial directions of the second and third constituents. The fifteen bivectors $f_af_b$ square to
$-\id$, are antisymmetric and close on $\so6$; each commutes with $X_1$, because the index sets are disjoint and $2\cdot4$ is even in
\eqref{eq:bladecomm}, so all fifteen preserve $P_1$ and act faithfully on its image. Let $\Omega_6=f_1f_4+f_2f_5+f_3f_6$ generate the pairing
$f_k\mapsto f_{k+3}$. Its centraliser in $\so6$ has dimension nine and is $\uu3$; the derived algebra has dimension eight and rank two, is
compact, closed and simple, hence $\su3$, and its eight generators act faithfully on the block.
\end{theorem}

\begin{theorem}[$P_1$ is vector-like undressed]
\label{thm:p1}
$X_1$ is a product of four generators, none of them interior, so $|I||J|-|I\cap J|=4$ is even for every interior generator and
$[X_1,\gamma_\mu]=0$. Hence $P_1$ is not axial and its matter is vector-like by \cref{prop:vectorlike}.
\end{theorem}

The obstruction is to the undressed action only.

\begin{theorem}[Colour with a chiral doublet on $P_1$]
\label{thm:p1chiral}
Keep the colour generators in the six negative directions of factors two and three, and in the three positive directions of factor
three set
\[
t_1=\tfrac12e^{(3)}_1e^{(3)}_2,\qquad t_2=\tfrac12e^{(3)}_1e^{(3)}_3,\qquad t_3=-\tfrac12e^{(3)}_2e^{(3)}_3 .
\]
With $P_{1,L}=P_1\tfrac12(\id-\Gamma)$ and $L_i=P_{1,L}t_i$,
\[
[L_i,L_j]=\epsilon_{ij}{}^kL_k,\qquad [C_A,L_i]=0,\qquad -\sum_iL_i^2=\tfrac34P_{1,L},\qquad L_iP_{1,R}=0
\]
for every colour generator $C_A$. Vector colour and a commuting, purely left-handed doublet action therefore coexist on $P_1$.
\end{theorem}

\begin{proof}
The $t_i$ are half time--time bivectors of factor three, so each squares to $-\tfrac14\id$ and they close on $\su2$. The colour generators
carry only negative indices of factors two and three, disjoint from $\{e^{(3)}_1,e^{(3)}_2,e^{(3)}_3\}$, and two even blades with disjoint
supports commute. $P_{1,L}$ is an idempotent of rank $128$ commuting with each $t_i$, so $L_i=P_{1,L}t_i$ inherits the bracket and
$-\sum L_i^2=P_{1,L}(-\sum t_i^2)=\tfrac34P_{1,L}$, the Casimir of a doublet. $P_{1,L}P_{1,R}=0$ gives the last identity.
\end{proof}

Ordering the family by $\dim\keff$ and recording whether the projector is axial,
\[
\Pfam\,(6)\ \text{axial},\quad P_{23}\,(9)\ \text{axial},\quad P_{12}\,(12)\ \text{not},\quad P_{13}\,(14)\ \text{axial},\quad P_1\,(24)\ \text{not},
\]
the two orderings are not reverses of one another, and what survives is a statement about these five projectors only: no axial one
carries $\su3$ in its effective bivector algebra, and $\su3$ appears on exactly one of them, the vector-like $P_1$. The pattern is
not a law of the arena. The next section exhibits a projector outside the family on which colour and chirality coexist.

% --- source: tex/s08_pstar#0 ((head))


% --- source: tex/s08_pstar#1 (A chiral Pati--Salam block)
\section{A chiral Pati--Salam block}
\label{sec:pstar}

The five projectors of \cref{sec:projectors} are one family among many. Enlarging the class to all commuting symmetric blade
involutions produces a projector on which colour and chirality coexist. The idea is to correlate spacetime chirality with an
\emph{internal} four-dimensional chirality, and to constrain the colour \emph{volume} rather than any colour direction.

% --- source: tex/s08_pstar#2 (Construction)
\subsection{Construction}

At three constituents take the interior gamma matrices $\gamma^\mu=(e^{(1)}_1,e^{(1)}_4,e^{(1)}_5,e^{(1)}_6)$ and split the remaining
fourteen directions into three groups:
\begin{center}
\begin{tabular}{@{}llcc@{}}
\toprule
group & directions & signs & generators $E_A$\\
\midrule
weak $W$ & $e^{(2)}_1,\ e^{(3)}_1,\ e^{(3)}_2,\ e^{(3)}_3$ & $++++$ & $7,13,14,15$\\
colour $C$ & $e^{(2)}_4,e^{(2)}_5,e^{(2)}_6,\ e^{(3)}_4,e^{(3)}_5,e^{(3)}_6$ & $------$ & $10,11,12,16,17,18$\\
auxiliary $A$ & $e^{(1)}_2,e^{(1)}_3,\ e^{(2)}_2,e^{(2)}_3$ & $++++$ & $2,3,8,9$\\
\bottomrule
\end{tabular}
\end{center}
With $J=\wI$ the ordered product of the fourteen internal generators and $\Gamma=J\gamma^0\gamma^1\gamma^2\gamma^3$, define the three
volumes
\[
\chi_W=e^{(2)}_1e^{(3)}_1e^{(3)}_2e^{(3)}_3,\qquad
\chi_C=e^{(2)}_4e^{(2)}_5e^{(2)}_6e^{(3)}_4e^{(3)}_5e^{(3)}_6,\qquad
q_A=e^{(1)}_2e^{(1)}_3 .
\]

\begin{figure}[t]
\centering
\begin{tikzpicture}[scale=0.95]
\foreach \k in {1,2,3}{\node at (1.3*\k,6.55) {\small $k=\k$};}
\node[anchor=south] at (2.6,6.85) {\small constituent};
\foreach \a in {1,...,6}{\node at (0.35,6.5-\a) {\small $e_{\a}$};}
% colour code: interior = gray, W = gcol, C = ocol, A = tcol
\def\pscell#1#2#3{\fill[#3] (1.3*#1-0.5,6.0-#2) rectangle (1.3*#1+0.5,6.0-#2+0.9);\draw[gray] (1.3*#1-0.5,6.0-#2) rectangle (1.3*#1+0.5,6.0-#2+0.9);}
\pscell{1}{1}{lgray}\pscell{1}{4}{lgray}\pscell{1}{5}{lgray}\pscell{1}{6}{lgray}
\pscell{1}{2}{tcol!35}\pscell{1}{3}{tcol!35}
\pscell{2}{1}{gcol!40}\pscell{2}{2}{tcol!35}\pscell{2}{3}{tcol!35}
\pscell{2}{4}{ocol!45}\pscell{2}{5}{ocol!45}\pscell{2}{6}{ocol!45}
\pscell{3}{1}{gcol!40}\pscell{3}{2}{gcol!40}\pscell{3}{3}{gcol!40}
\pscell{3}{4}{ocol!45}\pscell{3}{5}{ocol!45}\pscell{3}{6}{ocol!45}
\foreach \k in {1,2,3}{\foreach \a in {1,...,6}{\pgfmathtruncatemacro{\A}{6*(\k-1)+\a}\node at (1.3*\k,6.45-\a) {\scriptsize $E_{\A}$};}}
\draw[thick] (0.75,3.05) -- (4.45,3.05);
\node[anchor=west,align=left] at (5.0,5.6) {\small \textcolor{mgray}{\rule{8pt}{8pt}} interior $\gamma^\mu$: Lorentz};
\node[anchor=west,align=left] at (5.0,4.95) {\small \textcolor{gcol!70}{\rule{8pt}{8pt}} weak $W$: volume $\chi_W$, $\su2_L\oplus\su2_R$};
\node[anchor=west,align=left] at (5.0,4.3) {\small \textcolor{ocol!80}{\rule{8pt}{8pt}} colour $C$: volume $\chi_C$, $\so6\cong\su4$};
\node[anchor=west,align=left] at (5.0,3.65) {\small \textcolor{tcol!70}{\rule{8pt}{8pt}} auxiliary $A$: $q_A=e^{(1)}_2e^{(1)}_3$};
\node[anchor=west,align=left] at (5.0,2.6) {\small $\mathcal X=\Gamma\chi_W,\quad \mathcal Z=-J\chi_C,\quad \mathcal U=-Jq_A$};
\node[anchor=west,align=left] at (5.0,1.9) {\small $P_*=\tfrac18(\id+\mathcal X)(\id+\mathcal Z)(\id+\mathcal U)$, rank $64$};
\node[anchor=west,align=left] at (5.0,1.2) {\small temporal rows above the line, spatial below};
\end{tikzpicture}
\caption{The eighteen generators of $\Cl(9,9)$ arranged by constituent, and the grouping that defines the Pati--Salam block. The
interior four-plane carries the Lorentz action. The weak and colour groups enter $P_*$ only through their volumes, so every rotation
inside each of them survives the projection. The auxiliary group enters through the two-plane volume $q_A=e^{(1)}_2e^{(1)}_3$, not through the
volume of its four directions, and of its rotations only $E_2E_3$ and $E_8E_9$ survive ($E_2E_8$ anticommutes with $-Jq_A$, and
$P_*E_2E_8P_*=0$). The surviving bivectors, $15$ colour, $6$ weak, $E_2E_3$ and $E_8E_9$, span the $23$-dimensional algebra of
\cref{thm:PSalg} (\cref{thm:PS}).}
\label{fig:pstar}
\end{figure}

\begin{theorem}[A chiral projector carrying colour]
\label{thm:PS}
$\mathcal X=\Gamma\chi_W$, $\mathcal Z=-J\chi_C$ and $\mathcal U=-Jq_A$ are commuting symmetric traceless involutions, and all seven nonidentity products $\mathcal X$, $\mathcal Z$,
$\mathcal U$, $\mathcal{XZ}$, $\mathcal{XU}$, $\mathcal{ZU}$, $\mathcal{XZU}$ are traceless. Hence
\begin{equation}
\label{eq:Pstar}
P_*=\tfrac18(\id+\Gamma\chi_W)(\id-J\chi_C)(\id-Jq_A)
\end{equation}
is an orthogonal projector with $\Tr P_*=512/8=64$. It satisfies $[P_*,\Sigma^{\mu\nu}]=[P_*,J]=0$ for the six Lorentz generators
$\Sigma^{\mu\nu}=\tfrac12\gamma^\mu\gamma^\nu$, and $P_*\gamma^\mu P_*=0$, so it is axial; its two Weyl blocks have real rank $32$ each, and
$\chi_WP_*=\Gamma P_*$.
\end{theorem}

Rank is not automatic from the three factors being individually traceless; it is the products that decide, and here all of them are
traceless. The statement $\chi_WP_*=\Gamma P_*$ is the correlation of spacetime chirality with the internal chirality of the weak
directions.

% --- source: tex/s08_pstar#3 (The effective algebra)
\subsection{The effective algebra}

\begin{theorem}[The effective algebra and its Pati--Salam part]
\label{thm:PSalg}
All fifteen colour bivectors $e_ae_b$, $a,b\in C$, commute with $\mathcal X$, $\mathcal Z$ and $\mathcal U$, so the whole compact $\so6_C\cong\su4_C$ survives
and acts faithfully on the block; the colour volume $\chi_C$ is invariant under those rotations, so constraining it costs no colour
generator. With $w_1,\dots,w_4$ the weak directions, $\tau_i=-\tfrac12w_jw_k$ for cyclic $(i,j,k)$, $t^{L,R}_i=\tfrac12(\id\mp\chi_W)\tau_i$
and $T^{L,R}_i=P_*t^{L,R}_i$,
\[
[T^L_i,T^L_j]=\epsilon_{ijk}T^L_k,\qquad T^L_iP_R=0,\qquad -\sum_i(T^L_i)^2=\tfrac34P_L,
\]
and the mirror identities hold on the right. The $t^{L,R}_i$ are genuine bivectors, for example $t^L_3=-\tfrac14(w_1w_2+w_3w_4)$.
The full effective compact bivector algebra is
\[
\keff(P_*)=\su4\oplus\su2_L\oplus\su2_R\oplus\uu1^2,
\]
of dimension $23$, entirely compact, with a derived algebra of dimension $21$ and a two-dimensional centre. All $21$ simple generators
preserve the kinetic kernels $K^\mu=P_*\gamma^0\gamma^\mu P_*$, and $K^0=P_*$. (On the block of $P_*$ these are minus the compressed
one-constituent kernels of \cref{prop:kernels1}; the overall sign of the fermion Lagrangian is a convention, which changes neither the
field equation, nor the invariance condition $X^{\!\top}K^\mu+K^\mu X=0$, nor the one-momentum spectrum, see the discussion after
\cref{r1:prop:dirac}.)
\end{theorem}

\begin{proof}
Each colour bivector has even overlap with each of the three involution masks, so all fifteen commute with $\mathcal X,\mathcal Z,\mathcal U$ and descend
faithfully. For the weak part, $\chi_W^2=\id$ and $[\chi_W,\tau_i]=0$ make $\tfrac12(\id\mp\chi_W)$ idempotents commuting with the
$\tau_i$; the $\tau_i$ close on $\su2$ with $\tau_i^2=-\tfrac14\id$, so $-\sum_i(T^L_i)^2=\tfrac34P_L$, and $\chi_WP_*=\Gamma P_*$ places the
two ideals on opposite chiralities. The dimension $23$ is the full solution of $[B,P_*]=0$ over the $91$-dimensional space of internal
bivectors: the nullspace is $23$-dimensional, entirely compact, has trivial annihilator on the block, and has a derived algebra of
dimension $21$.
\end{proof}

\noindent Executable check: a10-1 (exact; \cref{app:methods}).

\begin{corollary}[One Pati--Salam family]
\label{cor:PSfamily}
The colour volume is not an involution: $\chi_C$ is a grade-six blade on six negative directions, so $\chi_C^2=-\id$. The
involution is
\[
\Gamma_C:=J\chi_C,\qquad \Gamma_C^2=\id,\qquad \Gamma_CP_*=-P_*,\quad\text{equivalently}\quad \chi_CP_*=JP_*.
\]
On the module with $J=+i$ each Weyl block then carries exactly four $\su4$ weights, all of colour chirality $\Gamma_C=-1$, each with
multiplicity $4=2\times2$, a weak doublet times a Weyl spinor, and the two blocks carry the same colour weights. The content is
\[
(\mathbf4,\mathbf2,\mathbf1)_L\oplus(\mathbf4,\mathbf1,\mathbf2)_R ,
\]
one Pati--Salam family including a right-handed neutrino \cite{PatiSalam}. Which of the two conjugate four-dimensional representations
is called $\mathbf4$ is a convention.
\end{corollary}

The decisive steps of \cref{thm:table,thm:PSalg} and of this corollary are statements that nullspaces have the stated dimensions. They are
certified exactly, in integer arithmetic: the $23$-generator basis of \cref{thm:PSalg} together with the completeness argument that every
blade commutes or anticommutes with each of the three involutions (a10-1); the centre and derived-algebra dimensions of every row of
\cref{thm:table} (a10-2); and the sign of the colour volume, rather than the quadratic Casimir, as what tells $\mathbf4$ from $\bar{\mathbf4}$
(a10-3).

The content is prior art. Embeddings of one chiral family together with the Lorentz group in a single orthogonal algebra are due to Maraner
\cite{Maraner04}, Nesti and Percacci \cite{NestiPercacci10} and Krasnov \cite[\S\S6.4--6.7]{Krasnov18}, and the same family is the fermion
content of the spectral Pati--Salam model \cite{ChamseddineConnesVanSuijlekom13}. What this realisation adds is the projector $P_*$ built from
three commuting involutions, with the normal charge on an extra positive four-plane, and the computations inside $V_*$; the comparison is
made in \cref{ch:ch14_related}.

% --- source: tex/s08_pstar#4 (Charges)
\subsection{Charges}

The compact algebra is realised by skew operators; charges are the corresponding symmetric operators $\hat t=-Jt$. On the block the
centre of $\keff(P_*)$ is spanned by the skew operators $JP_*$ and $J\Gamma P_*$, whose charges are $P_*$ and $\Gamma P_*$. Concretely
$q_A|_{V_*}=J$ and $E_8E_9|_{V_*}=J\Gamma$ on the image $V_*=\operatorname{im}P_*$.

\begin{proposition}[Colour intertwiner, $\BmL$, and the charge table]
\label{prop:PSledger}
$R=P_*\gamma^0e^{(2)}_1e^{(2)}_2$ satisfies $R^2=-P_*$, exchanges the two Weyl blocks and commutes with every colour generator, so the
colour representations of the two chiralities are equivalent and the content is $(\mathbf4,\mathbf2,\mathbf1)_L\oplus(\mathbf4,\mathbf1,
\mathbf2)_R$ rather than a $\mathbf4/\bar{\mathbf4}$ pairing. With
\[
H=-\bigl(e^{(2)}_4e^{(3)}_4+e^{(2)}_5e^{(3)}_5+e^{(2)}_6e^{(3)}_6\bigr),\qquad Q_c=-JH,
\]
the derived centraliser of $H$ inside colour $\so6$ is $\su3$, and $(Q_c-\id)(Q_c+3\id)P_*=0$: $Q_c$ has eigenvalues $+1$ and $-3$ on
the block, with real multiplicities $48$ and $16$, the triplet and the singlet. Setting $\BmL=Q_c/3$ and using the weak charges
$-JT^3_{L,R}$,
\[
Y=-J\,T^3_R+\tfrac12(\BmL),\qquad Q=-J\,T^3_L+Y ,
\]
gives one family:
\begin{center}
\begin{tabular}{@{}lccr@{}}
\toprule
field & $SU(3)\times SU(2)_L$ & $Y$ & internal Weyl components\\
\midrule
$Q_L$ & $(\mathbf3,\mathbf2)$ & $\tfrac16$ & $6$\\
$L_L$ & $(\mathbf1,\mathbf2)$ & $-\tfrac12$ & $2$\\
$u_R$ & $(\mathbf3,\mathbf1)$ & $\tfrac23$ & $3$\\
$d_R$ & $(\mathbf3,\mathbf1)$ & $-\tfrac13$ & $3$\\
$e_R$ & $(\mathbf1,\mathbf1)$ & $-1$ & $1$\\
$\nu_R$ & $(\mathbf1,\mathbf1)$ & $0$ & $1$\\
\bottomrule
\end{tabular}
\end{center}
Four real spinor coordinates make one internal Weyl component, so rank $64$ is sixteen internal Weyl components, eight left and eight
right, and not sixty-four particles.
\end{proposition}

The $\su3$ embedding, the right-handed Cartan direction and the hypercharge combination are chosen; what is established is their
mutual compatibility. In the complex structure $J$ every one of the sixteen Weyl fermions has charge one: $J$ is fermion number,
$F=3B+L$, and the baryon number appears as $4B=F+(\BmL)$.

\begin{proposition}[The two central factors cannot be gauged]
\label{prop:PSanom}
For a central charge $a+b\Gamma$, corresponding to the skew element $aJ+bJ\Gamma$, the anomaly coefficients on the block, per Weyl
fermion with every field written left-handed, are
\[
\begin{gathered}
SU(2)_L^2:\ 2(a-b),\qquad SU(2)_R^2:\ -2(a+b),\qquad SU(4)^2:\ -2b,\\
\text{cubic}:\ -16(3a^2b+b^3),\qquad \text{gravitational}:\ -16b ,
\end{gathered}
\]
and the mixed anomalies with two central and one Pati--Salam generator vanish. All vanish only for $a=b=0$. The Pati--Salam algebra itself
is anomaly free on the block: the $SU(4)^3$ anomaly cancels between two fundamentals and two antifundamentals in the all-left convention,
and each weak $SU(2)$ sees four doublets, so the global anomaly is absent \cite{Witten82}. Hence, with the block as the matter content and
within the $23$-dimensional effective algebra of \cref{thm:PSalg}, $\su4\oplus\su2_L\oplus\su2_R$ is the largest anomaly-free gauge algebra
containing Pati--Salam, and the two central directions remain global symmetries. This is not a classification of all algebras acting on the
block.
\end{proposition}

The obstruction is stated for this matter content; additional matter or another cancellation mechanism is not excluded. The two central
directions, fermion number $J$ and the chiral charge $J\Gamma$, are classical global symmetries with mixed anomalies, not exact quantum
symmetries. The fermion-number current has a triangle anomaly with $SU(2)_L$ and $SU(2)_R$ already in perturbation theory \cite{Adler69}; what
holds to all orders of perturbation theory is the corresponding charge selection rule in the trivial topological sector, which instantons
violate \cite{tHooft76}.

\begin{proposition}[Two rank-64 projectors, two orbits]
\label{prop:orbits}
Under the compact group of the connection sector, whose Lie algebra is the maximal compact $\so8\oplus\so6$ of dimension $43$, the orbit
through $\Pfam$ has dimension $43-9=34$ and the orbit through $P_*$ has dimension $43-23=20$. Equal rank does not put two projectors on
one orbit.
\end{proposition}

\begin{remark}[What the construction settles]
\label{rem:PSsettles}
It settles an existence question: colour on a chiral block is not obstructed, and no counting argument based on the five-projector family
can be a theorem. It does not select $P_*$; it does not derive the hypercharge embedding, which is compatible with anomaly cancellation but
chosen; it gives one family and not three; and by \cref{cor:nopredict} exhibiting a known representation content is a consistency result,
not evidence. One further fact will matter in \cref{sec:matter}: $\gamma^0$ commutes with all $91$ internal bivectors and with $J$ and the
rotations, and anticommutes with $\Gamma$ and the boosts. It maps $V_*$ onto an orthogonal rank-$64$ subspace on which the same internal
algebra acts with $\Gamma$ reversed. The chirality of the block is put in by the choice of $P_*$.
\end{remark}

% --- source: tex/s09_charges#0 ((head))


% --- source: tex/s09_charges#1 (Charges and anomalies)
\section{Charges and anomalies}
\label{sec:charges}

The weak and colour algebras of a block are fixed by the projector; the abelian charges are not. This section asks how much the
algebra says about them, first inside the arena and then in a layer that uses nothing from the arena at all.

% --- source: tex/s09_charges#2 (A pairing obstruction)
\subsection{A pairing obstruction}
\label{sec:pairing}

The family block of $\Pfam$ carries three abelian charges. The centre of $\keff(\Pfam)$ is spanned by the block images of
$X_a\in\{\tfrac12J^{(1)},\tfrac12J^{(3)},\tfrac12e^{(3)}_5e^{(3)}_6\}$, and their charges $q_a:=-J(2X_a)\Pfam$ are commuting symmetric
involutions of the block, with eigenvalues $\pm1$ and $\Tr q_a=0$. The cubic anomaly of $u=aq_1+bq_2+cq_3$ is the chiral trace
\[
\mathcal A(u^3)=\tfrac14\bigl[\Tr(u^3P_L)-\Tr(u^3P_R)\bigr],\qquad P_{L,R}=\Pfam\tfrac12(\id\mp\Gamma),
\]
per Weyl component, right-handed components entering with conjugated charge when every field is written left-handed. Since
$q_a^2=\Pfam$ and the linear chiral traces vanish, only the mixed term survives:
\[
\mathcal A(u^3)=6abc\,\mathcal A(q_1q_2q_3)=96\,abc ,
\]
while the plain trace $\Tr(u^3)$ vanishes for every $u$ and carries no information about anomalies. The anomaly-free locus is the union
of the three coordinate planes $abc=0$, and on it there is nothing to select.

\begin{theorem}[Conjugate pairing]
\label{thm:pairing}
Let a module carry three commuting involutions $q_a$ with eigenvalues $\pm1$ and $\Tr q_a=0$. For every $u=aq_1+bq_2+cq_3$ with
$abc=0$ the weight multiset of $u$ is symmetric, $\mathrm{mult}(w)=\mathrm{mult}(-w)$.
\end{theorem}

\begin{proof}
Take $c=0$; only the joint multiplicities $M_{s_1s_2}$ of $(q_1,q_2)$ matter. The trace conditions give
$M_{++}+M_{+-}-M_{-+}-M_{--}=0$ and $M_{++}-M_{+-}+M_{-+}-M_{--}=0$; adding and subtracting, $M_{++}=M_{--}$ and $M_{+-}=M_{-+}$. The
weight of $(s_1,s_2)$ is $as_1+bs_2$, and $(s_1,s_2)\mapsto(-s_1,-s_2)$ negates it while preserving multiplicities.
\end{proof}

The proof uses only the three linear traces, with no reference to rank, rung or realisation, so within its hypotheses no
enlargement of the block and no ascent of the tower repairs it. Charge pairing is all that is concluded: calling a full matter
representation vector-like would also require the chirality convention and the other gauge factors. And the hypotheses are not
decorative.

\begin{example}[Vanishing traces do not imply pairing]
\label{ex:multiset}
The multiset $\{(-3)^{\times5},1^{\times10},5^{\times1}\}$ has $\sum q=-15+10+5=0$ and $\sum q^3=-135+10+125=0$, yet it is not symmetric under
$q\mapsto-q$. It is realised by three commuting involutions with joint signs and multiplicities $(+,+,+){:}\,1$, $(-,-,+){:}\,10$,
$(+,-,-){:}\,5$: these have traces $(\Tr q_1,\Tr q_2,\Tr q_3)=(-4,-14,6)$, so the hypotheses of \cref{thm:pairing} fail, while
$Y=q_1+q_2+3q_3$ has spectrum $\{5^{\times1},1^{\times10},(-3)^{\times5}\}$ with $\Tr Y=\Tr Y^3=0$. Leaving the hypotheses therefore need not
violate the gravitational condition, which constrains $\Tr Y$ for the gauged combination and not $\Tr q_a$ for each basis charge
\cite{AlvarezGaumeWitten,Bilal}.
\end{example}

% --- source: tex/s09_charges#3 (Dressing, and a benchmark)
\subsection{Dressing, and a benchmark}
\label{sec:benchmark}

The escape from \cref{thm:pairing} is to leave its hypotheses. On the family block let
\[
t^1=-J\cdot\tfrac12e^{(2)}_5e^{(2)}_6,\qquad t^2=-J\cdot\tfrac12e^{(2)}_6e^{(2)}_4,\qquad t^3=-J\cdot\tfrac12e^{(2)}_4e^{(2)}_5
\]
be the weak charges, whose skew counterparts, the rotations of the spatial directions of constituent two, span the derived $\su2$ of
$\keff(\Pfam)$; let $\varepsilon=\pm1$ be an orientation sign, and put $t^a_{L,R}=P_{L,R}t^a$,
\[
B_0=\frac\varepsilon6\,\Gamma(q_1+q_2+q_3),\qquad Y=B_0+t^3_R,\qquad Q_{\rm em}=t^3_L+Y .
\]
For the displayed charge table we choose $\varepsilon=+1$ in the stated orientation; reversing $\varepsilon$ reverses the $B_0$ charges, and the
conventions must then be changed before the Standard-Model identification below is reused. The cell sums $\varepsilon\sum_aq_a$ are then $\{-1,-1,-1,3\}$ on the $\Gamma=-1$ half and $\{-3,1,1,1\}$ on the other, which gives
$B_0|_{\Gamma=-1}=\{\tfrac16,\tfrac16,\tfrac16,-\tfrac12\}$ and, after adding $t^3_R=\pm\tfrac12$ on the right-handed cells,
$Y\in\{0,-1,\tfrac23,-\tfrac13\}$ there. Over the sixteen left-handed Weyl components $\Tr Y=\Tr Y^3=\Tr T_3^2Y=0$ with four doublets, so all
four anomaly conditions hold, and $\Tr T_3^2/\Tr Q^2=3/8$. This hypercharge is adopted, not derived.

\begin{proposition}[The escape]
\label{prop:escape}
\Cref{thm:pairing} constrains undressed elements of $\uu1^3$, which lie in the vector class. The charge $B_0$ carries a factor $\Gamma$ and is
axial, not an element of $\uu1^3$, so the cubic-trace formula does not apply to it. Once chirality-dressed endomorphisms are admitted,
which abelian charges can be gauged changes, and no-go results proved in the undressed grade-two setting do not carry across; a well-posed
search must say which Clifford grades are permitted alongside $Y^\dagger=Y$, $[Y,t_{La}]=0$ and $[Y,C_A]=0$.
\end{proposition}

Since $P_L+P_R=\id$, the chirality projectors on the weak axis cancel in $Q_{\rm em}=t^3+B_0$ when the same axis is used in both terms. The
dressing of $B_0=\tfrac\varepsilon6\Gamma(q_1+q_2+q_3)$ does not cancel, so $Q_{\rm em}$ still contains $\Gamma$. The value $3/8$ is arithmetic on a
charge table; it becomes $\sin^2\theta_W$ only after the equality of the normalised couplings is imposed at a matching scale, as a boundary
condition, which this chapter does not supply (the adopted action of \cref{ch:ch07_vacuum} supplies it by choice, \cref{thm:compat}). Five choices enter the
benchmark: the projector, the dressing, which abelian direction is $Y$, which right-handed axis, and the overall sign.

% --- source: tex/s09_charges#4 (A model-independent layer)
\subsection{A model-independent layer}
\label{sec:modelindep}

The following uses nothing from the arena. It calibrates how much of the hypercharge problem is generic.

\begin{theorem}[The perturbative hypercharge-anomaly variety of one family]
\label{thm:variety}
For one family with $N_c\ge2$ colours and a right-handed neutrino (for $N_c=1$ there is no colour mixed condition), with left-handed
hypercharges $q,l$ and right-handed $y_u,y_d,y_\nu,y_e$,
the three linear conditions give $l=-N_cq$, $y_d=2q-y_u$, $y_e=2l-y_\nu$. With $h_q=y_u-q$ and $h_\ell=y_\nu-l$ the cubic condition
factorises,
\[
\mathcal A_{Y^3}=6N_cq\,(h_\ell^2-h_q^2),
\]
using $2q^3-y_u^3-(2q-y_u)^3=-6qh_q^2$. The solution set is the union of three planes: $\Pi_A$ ($h_\ell=h_q$), $\Pi_B$ ($h_\ell=-h_q$) and
$\Pi_C$ ($q=l=0$). On $\Pi_A$ the four Yukawa relations $y_u=q+h$, $y_d=q-h$, $y_\nu=l+h$, $y_e=l-h$ hold, so they are consequences of anomaly
cancellation rather than assumptions; $\Pi_B$ is the image of $\Pi_A$ under the exchange $\nu_R\leftrightarrow e_R$. At $N_c=3$ and
$h=\tfrac12$, $\Pi_A$ is the line $Y=Y_{\rm SM}+\lambda(\BmL)$.
\end{theorem}

Anomaly cancellation does not supply the scalar, its hypercharge, or nonvanishing couplings. Quark Yukawa couplings alone do not exclude
$\Pi_C$ either: at $N_c=3$ the point $(q,l,y_u,y_d,y_\nu,y_e)=(0,0,\tfrac12,-\tfrac12,1,-1)$ with scalar hypercharge $\tfrac12$ cancels all four
anomalies and permits both quark Yukawas, yet $h_q=\tfrac12\neq h_\ell=1$. What excludes it is integrality: with $q=0$ the left-handed quarks
carry $\pm\tfrac12$ and a three-quark singlet has half-integral charge.

\begin{theorem}[Without $\nu_R$ hypercharge is quantised on the generic branch]
\label{thm:nonu}
Deleting $\nu_R$ gives $y_e=2l$ and $\mathcal A_{Y^3}=6N_cq(N_c^2q^2-h_q^2)$, so on the generic branch $h_q=\pm N_cq$ is locked and
hypercharge is fixed up to scale and relabelling: at $N_c=3$, $q=\tfrac16$ gives the Standard-Model values with no residual $\BmL$ direction.
\end{theorem}

This is the classical charge-quantisation result \cite{FootLewVolkas,GengMarshak}. Its bearing here is that the difficulty of selecting
hypercharge on a family block exists because the block carries sixteen Weyl components rather than fifteen: the sterile neutrino is both
the block's virtue and the source of the ambiguity.

\begin{proposition}[Integrality, and what it does not fix]
\label{prop:integrality}
On $\Pi_A$ with $Q_{\rm em}=T^3_L+Y$, the leptons carry $Q_\nu=-\lambda$ and $Q_e=-1-\lambda$, so integrality of the charged lepton forces
$\lambda\in\Z$, and colour triality then gives integral charge to every colour singlet. It does not fix the fractional part of the quark
charges, since $Q_u=\tfrac23+\lambda/3$, and neutrality does not close the argument for composites: $Q(ddd)=\lambda-1$ and $Q(uuu)=\lambda+2$, so
neutral colour-singlet operators of odd fermion parity exist at $\lambda=1$ and $\lambda=-2$. A neutral \emph{elementary} colour-singlet
fermion exists if and only if $\lambda\in\{0,-1\}$. These give $(Q_u,Q_d,Q_\nu,Q_e)=(\tfrac23,-\tfrac13,0,-1)$ and $(\tfrac13,-\tfrac23,1,0)$;
exchanging $\nu$ and $e$ alone does not map one to the other, but charge conjugation together with $u\leftrightarrow d$ and
$\nu\leftrightarrow e$ does.
\end{proposition}

Three other routes fail to single out the Standard Model. The kernel of the representation is not a selection principle: its order
depends on $\lambda$ through the integer normalisation of the weights. Nor is a $\Z_6$ kernel evidence for three colours: with a neutral
$\nu_R$ the kernel is $\Z_{2N_c}$ for odd $N_c$ and $\Z_{N_c}$ for even $N_c$. And integrality does not select three either: on the standard
branch every $N_c$-quark colour singlet has integral charge exactly for odd $N_c$, and the number $N_c+1$ of weak doublets per family, which
the global $SU(2)$ anomaly requires to be even \cite{Witten82}, is even exactly for odd $N_c$. All three facts hold at every odd $N_c$. With
$N_g$ generations the parameter count depends on how the family is specified. If anomalies cancel only after summing over generations, the
perturbative anomaly variety of unnormalised hypercharges has local dimension $6N_g-4$ at the Standard-Model point ($14$ for three
generations), before an overall normalisation is fixed. A $2N_g$-dimensional family is obtained
with one common Yukawa difference $h$, $y_u=q_i+h$, $y_d=q_i-h$, $y_\nu=l_i+h$, $y_e=l_i-h$ with $\sum_i(N_cq_i+l_i)=0$, and it contains $L_e-L_\mu$;
imposing the one-family branch generation by generation also gives $2N_g$ parameters but excludes $L_e-L_\mu$. The anomaly-free directions
include $L_e-L_\mu$ and $L_\mu-L_\tau$ \cite{HeJoshiLewVolkas}, so family universality is an input that grows with $N_g$. The global structure of the gauge group is a separate
datum \cite{Tong,GaiottoKapustinSeibergWillett}.

% --- source: tex/s09_charges#5 (Charge at one constituent)
\subsection{Charge at one constituent}

Complexifying, the spinor module of one constituent is $\Lambda^\bullet\C^3$, and its even half is spanned by
$|0\rangle,|12\rangle,|13\rangle,|23\rangle$. The diagonal Fock charge $N_1-N_2$ has spectrum $\{0,0,+1,-1\}$ there, the charge skeleton of
$(\nu,\bar\nu,e^+,e^-)$.

\begin{theorem}[No unique charge at one constituent]
\label{thm:fock}
The condition $\spec Q=\{0,0,+1,-1\}$ on the even half, for $Q=\sum_iq_iN_i+c$, has exactly twelve solutions, forming two orbits of
$S_3\times\{\pm1\}$: the six differences $Q=N_i-N_j$, with $c=0$, and the six signed sums $Q=\pm(N_i+N_j-1)$, with $c=\mp1$. Hodge
conjugation $C$ maps exterior degree $k$ to $3-k$ and so exchanges the even and odd halves; the relation $CQC^{-1}=-Q$, equivalent to
$\sum_iq_i+2c=0$, holds on the whole exterior algebra for both orbits, not as an operation inside the even four-state sector, and does not
separate them.
\end{theorem}

\begin{proof}
Summing the four values gives $4c+2\sum_iq_i=0$. If the neutral pair is $\{|0\rangle,|ij\rangle\}$ then $c=0$ and $q_i+q_j=0$, which forces
the third coefficient to vanish; if it is $\{|ij\rangle,|ik\rangle\}$ then $q_j=q_k$, $c=-q_j$ and $q_i=0$.
\end{proof}

The charge skeleton alone does not identify neutrinos and electrons. The further step, $Q|0\rangle=0$ and hence $c=0$, is the separate
assumption that the state identified with the neutrino is the Fock vacuum.
Choosing between the orbits is an identification, not a derivation, and it has the same shape as the hypercharge problem two rungs up,
differing in size and not in kind.


\begin{remark}[Relation to Krasnov's construction]
\label{rem:krasnovpointer}
Krasnov describes one generation, with the right-handed neutrino, as a real Weyl spinor of $\mathrm{SO}(7,7)$ realised on the even forms
on $\R^7$; the Pati--Salam group appears as the compact part of the commutant $\mathrm{SO}(4,6)$ of the Lorentz group \cite{Krasnov18}.
The block $V_*$ has the same real dimension, $64$, but it is not a module of $\mathfrak{so}(4,6)$: of the $45$ bivectors of
$\mathfrak{so}(4,6)$ on the weak and colour directions, only the $21$ compact ones preserve $P_*$, and each of the $24$ non-compact ones
maps the block to its complement (check z3-9). As representations of the common algebra $\so{1,3}\oplus\so4\oplus\so6$ the two modules are
nevertheless equivalent: both are the underlying real representation (the realification) of $(\mathbf2,\mathbf1)\otimes[(\mathbf4,\mathbf2,\mathbf1)\oplus(\bar{\mathbf4},\mathbf1,\mathbf2)]$ (\cref{cor:PSfamily}),
so an intertwiner exists by complete reducibility. It has not been written down, and whether it carries the kinetic structures is open; the
non-compact generators of $\so{4,6}$ obstruct only its extension to the ambient action. The comparison is made in \cref{ch:ch14_related}.
\end{remark}

